The difference does not show that either percentage is the unique ``true'' pass rate; it shows that the construct changes when adverse statuses are translated into simulated grades. For Paper 3, the primary distributional estimand is therefore the distribution of awarded numeric grades, while adverse-status incidence becomes a parallel classroom characteristic. The augmented outcome remains valuable as a sensitivity analysis because a classroom with many withdrawals may look unusually successful if only completed numeric grades are inspected.

\subsection{Exploratory distributional clusters}

The historical report computed pairwise KS distances between professor grade distributions and used K-Medoids to search for groups of similar distributional profiles. Its reported Silhouette Score was maximal at $k=3$, and the corresponding plots showed three sets of distributions that were more similar within groups than across the full collection. This result is evidence that heterogeneity was not visually random, but the original clustering combined professor records across academic contexts and should therefore be interpreted only as a hypothesis-generating result.

A confirmatory cluster analysis, if retained in the final paper, will be secondary to direct distance summaries and will operate within explicitly comparable strata. The paper does not require a discrete cluster solution to establish contextual heterogeneity; continuous variation in classroom distributions is itself substantively meaningful.

\subsection{Period belongs in the grading context}

The historical report also follows the mean grade of the same instructor across periods and shows visible temporal movement. Although an illustrative series cannot estimate population-level stability, it demonstrates why the analysis should not define grading context by instructor alone. If an instructor's distribution moves as course cohort, assessment design or period conditions change, pooling that instructor's records across years can create a reference distribution that no individual classroom actually experienced.

\draftnote{The controlled rerun should replace the illustrative argument with instructor-course combinations observed in at least three periods, reporting between-period ranges and distribution-distance summaries, along with a stability coefficient or variance-component estimate when estimable.}

\subsection{How much does standardisation change student classification?}

The historical project used classroom-standardised grades because raw grades appeared insufficiently comparable across local contexts, but it did not quantify the measurement consequence of that choice directly. Paper 3 will therefore report how many observations change pooled percentile rank or quintile/decile classification when raw grade is replaced by classroom-relative $Z$. This is the most direct link between the heterogeneity diagnosis and downstream educational analysis: if reclassification is negligible, contextual differences may have little practical consequence for many uses; if it is substantial, raw-grade pooling can alter which students appear relatively high or low performing.

\draftnote{Is reclassification concentrated among observations near institutional thresholds such as 7.5, or is it distributed across the grade range? A threshold-concentrated result would have different implications from broad rank movement and should be reported separately.}

\section{Discussion}

The retained evidence supports a measurement problem that is narrower than a claim about instructor quality but more consequential than a cosmetic difference in grade distributions. Undergraduate mathematics grades are recorded on a common numerical scale, yet the historical data show that the distributions producing those numbers differ across instructor contexts and can change over time. Assessment scholarship provides several plausible reasons for this pattern, including variation in criteria, holistic academic judgement and local assessment cultures \citep{bloxham2011,bloxham2016,lipnevich2020,ylonen2018}. The administrative data cannot distinguish those mechanisms from student sorting or differences in realized achievement, so the defensible empirical object is context-associated grade variation rather than instructor severity.

This distinction matters because educational research frequently uses grades as outcomes without modelling the environments in which they were produced. A regression on raw grades can therefore combine at least two sources of variation: differences in student performance and differences associated with classroom assessment context. Restricting comparisons to course-period strata, including classroom effects or using classroom-relative standardisation can reduce the second component, but each strategy changes the target of inference. In particular, $Z_{ic}$ answers where a student sits within the local classroom distribution, not what absolute level of mathematical knowledge the student possesses.

The difference between reliability and comparability is also important. High reliability of accumulated college grades \citep{beatty2015} is compatible with local grading heterogeneity because a stable aggregate can be built from components whose contexts differ systematically. Conversely, finding heterogeneous classroom distributions does not imply that grades are unreliable for progression decisions or useless as institutional records. Paper 3 instead asks whether raw course grades are exchangeable enough for analytical comparisons that ignore instructor-course-period context.

For assessment practice, the findings could motivate attention to coordination, moderation and shared standards, but the paper should not infer that uniform distributions are a desirable policy target. Mechanical equalisation of means could conceal genuine cohort differences or distort criterion-referenced assessment. Research on grading variation suggests that moderation and shared interpretation can reduce some inconsistency while recognising the role of professional judgement \citep{hunter2011,bloxham2016}. The more immediate institutional implication is analytical transparency: whenever course grades are pooled, researchers should state whether and how grading context is handled and what interpretation remains after adjustment.

\section{Limitations and remaining analyses}

Several limitations are structural. First, students are not randomly allocated to instructors, so between-classroom distributional differences cannot be interpreted as causal grading effects. Second, the administrative data record final outcomes but not the full assessment design, weights of examinations and assignments, marking rubrics, moderation processes or changes in course content; those unobserved features are part of the mechanism rather than statistical noise. Third, classroom standardisation removes location and scale by construction and can hide absolute differences in mastery if those differences are real. Fourth, small classrooms yield unstable distribution estimates, making the minimum-size sensitivity essential. Fifth, nonnumeric adverse outcomes are substantively informative but do not belong naturally on the same measurement scale as awarded numeric grades, so any augmented outcome requires explicit modelling assumptions.

